\chapter{Introducción y Objetivos}
\label{chap:introduccion}

El diseño de etapas analógicas amplificadoras con transistores de efecto de campo metal-óxido-semiconductor (MOSFET)
constituye un pilar esencial en la electrónica discreta e integrada moderna. A diferencia de los transistores
bipolares de juntura (BJT), controlados esencialmente por corriente de base, los transistores MOSFET son dispositivos
controlados por tensión, cuya compuerta se encuentra aislada eléctricamente del canal por una delgada capa dieléctrica
de dióxido de silicio ($\mathrm{SiO}_2$). Esta propiedad física fundamental dota al dispositivo de una impedancia de
entrada de corriente continua prácticamente infinita ($I_G \approx 0$).

Entre las topologías fundamentales de amplificación con MOSFET, la configuración en Surtidor Común (CS, por sus siglas
en inglés \textit{Common-Source}) es la más ampliamente difundida para proveer una sustancial ganancia de tensión con
inversión de fase ($180^\circ$) respecto a la señal de excitación, resultando análoga a la configuración en emisor común
de los transistores bipolares.

\section{Objetivos del Trabajo Práctico}

El presente trabajo práctico de laboratorio plantea como meta el diseño, modelado, simulación y contrastación
experimental completa de una etapa amplificadora en surtidor común basada en el transistor MOSFET de canal N tipo
enriquecimiento 2N7000. Los objetivos específicos fijados son:

\begin{enumerate}
  \item Relevar experimentalmente en laboratorio la curva característica de transferencia en zona de saturación del
    transistor 2N7000, deduciendo mediante regresión lineal la tensión umbral de compuerta ($V_{th}$) y el parámetro de
    transconductancia ($K$).
  \item Dimensionar analíticamente la red de polarización en corriente continua mediante divisor resistivo de compuerta
    y resistencia de autorpolarización de surtidor ($R_S$), garantizando estabilidad térmica y fijando el punto estático
    de reposo.
  \item Normalizar los valores resistivos obtenidos a componentes comerciales disponibles de la serie estándar E12,
    evaluando analíticamente los desvíos inducidos sobre el punto de reposo $Q(V_{DSQ}, I_{DQ})$.
  \item Simular el circuito en LTspice tanto en punto de operación estático (.op) como en régimen transitorio (.tran) ante
    excitaciones senoidales de $1\kHz$, validando el funcionamiento en zona de saturación y la ganancia dinámica.
  \item Trazar y analizar las rectas de carga estática de corriente continua y dinámica de corriente alterna,
    estableciendo las cotas de excursión lineal simétrica y verificando la no incursión en zona de triodo ni en corte.
  \item Desarrollar el modelo lineal de pequeña señal a frecuencias medias, derivando analíticamente las expresiones de
    la ganancia de tensión ($A_v$), ganancia de corriente ($A_i$), impedancia de entrada ($Z_i$) e impedancia de salida
    ($Z_o$).
  \item Implementar físicamente el circuito en el laboratorio sobre placa de ensayo y relevar con instrumental de banco
    las tensiones nodales continuas y las amplitudes de alterna.
  \item Realizar una contrastación cuantitativa entre la teoría analítica, las simulaciones computarizadas en
    LTspice y las mediciones de laboratorio.
\end{enumerate}
